\chapter{Introducción}
\label{chap:introduction}

\section{Contexto}

La contratación pública constituye una parte relevante de la actividad
económica española. En 2024 representó el 10,92\,\% del producto interior bruto
y se licitaron 206.242 contratos por un presupuesto base, sin impuestos, de
113.091,40 millones de euros \cite{oirescon2025}. Estas cifras permiten
dimensionar el dominio en el que se enmarca LIA, pero no implican que todas las
licitaciones presenten la misma complejidad ni que sean oportunidades adecuadas
para cualquier empresa.

La participación en un procedimiento requiere interpretar la documentación del
expediente. La Ley 9/2017 de Contratos del Sector Público establece que los
criterios y requisitos mínimos de solvencia económica y financiera, así como de
solvencia técnica o profesional, deben indicarse en el anuncio o la invitación y
detallarse en los pliegos. También exige especificar los medios admitidos y los
umbrales aplicables para acreditarlos \cite{lcsp2017}. Por tanto, cualificar una
licitación no consiste únicamente en localizar una cifra: exige relacionar los
requisitos concretos del contrato con la información disponible sobre la
empresa licitadora.

Además de la solvencia, el análisis debe considerar el objeto del contrato, sus
condiciones de ejecución, los importes y plazos relevantes y los criterios de
adjudicación. Esta información puede encontrarse distribuida entre anuncios,
pliegos administrativos, prescripciones técnicas y documentación
complementaria. La revisión conjunta de esas fuentes permite decidir si resulta
razonable preparar una oferta y qué aspectos requieren un estudio posterior.

La cualificación de una licitación es, por tanto, una actividad intensiva en
conocimiento. Los equipos responsables deben localizar información distribuida
entre diferentes documentos y relacionarla con las características de la
empresa que valora presentarse. Cuando el proceso se lleva a cabo de forma
manual, requiere una dedicación considerable y queda expuesto a posibles
errores de omisión.

En este contexto surge LIA, una plataforma web de gestión de oportunidades
comerciales centrada en licitaciones públicas. El proyecto combina la
organización de oportunidades mediante pipelines configurables y flujos de
trabajo estructurados con el uso de inteligencia artificial generativa para
apoyar las tareas de mayor carga cognitiva del proceso de análisis.

En esta memoria, una \emph{oportunidad} es una licitación registrada por el
equipo, y la \emph{cartera} es el conjunto de oportunidades que gestiona una
cuenta. El \emph{pipeline} organiza sus estados comerciales, mientras que el
\emph{workflow} define los pasos y acciones de su análisis. La
\emph{cualificación} consiste en estudiar la documentación y conservar
resultados revisables que apoyen la decisión de presentarse, sin sustituir el
juicio de la persona experta.

Los modelos generativos pueden producir respuestas plausibles que sean
incorrectas o inconsistentes, un riesgo que el NIST denomina
\emph{confabulación} \cite{autio2024}. En consecuencia, LIA plantea la IA como
un mecanismo de apoyo integrado en un proceso controlado: recibe el contexto
documental seleccionado, produce resultados estructurados, registra sus estados
de ejecución y permite que el usuario revise, edite o reintente el resultado.
La decisión final sobre la presentación a una licitación permanece fuera del
ámbito de automatización del sistema.

\begin{figure}[htbp]
  \centering
  \resizebox{\textwidth}{!}{
    \begin{tikzpicture}[
  node distance=0.46cm,
  every node/.style={font=\small\sffamily},
  flowbox/.style={
    draw=TFGBlue,
    rounded corners=2pt,
    fill=TFGLightBlue,
    minimum height=1.72cm,
    text width=2.58cm,
    align=center,
    inner sep=5pt
  },
  result/.style={
    flowbox,
    fill=TFGBlue,
    text=white
  },
  flow/.style={
    -{Latex[length=2.5mm]},
    semithick,
    draw=TFGAccent
  }
]
  \node[flowbox] (documents) {
    \textbf{Carga de documentación}\\[2pt]
    Pliegos, anuncios y memorias
  };
  \node[flowbox, right=of documents] (extraction) {
    \textbf{Extracción}\\[2pt]
    Identificación de información clave
  };
  \node[flowbox, right=of extraction] (assessment) {
    \textbf{Evaluación}\\[2pt]
    Solvencia y viabilidad preliminar
  };
  \node[flowbox, right=of assessment] (criteria) {
    \textbf{Análisis}\\[2pt]
    Requisitos y criterios de adjudicación
  };
  \node[result, right=of criteria] (decision) {
    \textbf{Resultado}\\[2pt]
    Documento de apoyo\\[-1pt]
    para su revisión
  };

  \draw[flow] (documents) -- (extraction);
  \draw[flow] (extraction) -- (assessment);
  \draw[flow] (assessment) -- (criteria);
  \draw[flow] (criteria) -- (decision);
\end{tikzpicture}

  }
  \caption{Flujo funcional propuesto para el análisis de una licitación.}
  \label{fig:tender-analysis-flow}
\end{figure}

La \Cref{fig:tender-analysis-flow} resume el flujo funcional descrito en la
propuesta registrada del proyecto. Los bloques representan actividades de
análisis y no decisiones completamente autónomas: el resultado es un documento
de apoyo que debe ser revisado por una persona experta.

\section{Motivación}

La motivación principal del proyecto parte de la necesidad de reducir el
esfuerzo dedicado a revisar y cualificar licitaciones públicas. Cuando la
información relevante está repartida entre varios documentos, el equipo debe
localizarla, contrastarla y trasladarla a un criterio de decisión. Estructurar
esta actividad reduce la dependencia de revisiones informales y permite
conservar el estado y el resultado de cada paso del análisis.

La inteligencia artificial generativa ofrece un mecanismo para producir
respuestas a partir de instrucciones y documentación de contexto. Sin embargo,
la posibilidad de obtener resultados incorrectos impide tratarla como una
fuente infalible. Por este motivo, LIA no se plantea como una conversación
aislada con un modelo, sino como una plataforma que relaciona cada generación
con una oportunidad, una acción del workflow, unos documentos seleccionados y
un estado verificable.

Otro aspecto que motiva el proyecto es la necesidad de equilibrar la capacidad
de los proveedores de inteligencia artificial generativa con el coste y la
latencia asociados a su utilización. Este objetivo requiere una evaluación
reproducible con los mismos documentos y criterios. La comparación experimental
del \Cref{chap:testing} y su análisis por operación en el
\Cref{sec:ai-model-selection} fundamentan la selección adoptada, sin confundir
una menor tarifa con una mayor fiabilidad.

\section{Objetivos del proyecto}

El objetivo general del presente Trabajo Fin de Grado es aligerar el proceso de
análisis y toma de requisitos de licitaciones públicas mediante el desarrollo
de una plataforma web de gestión de oportunidades comerciales asistida por
inteligencia artificial generativa.

Para alcanzar este objetivo general se establecen los siguientes objetivos
específicos:

\begin{itemize}
  \item Automatizar la detección de solvencia económica necesaria para valorar
        la presentación de una empresa a una licitación.
  \item Optimizar la relación entre el coste y el rendimiento derivado del uso
        de proveedores de inteligencia artificial generativa.
  \item Facilitar la creación de flujos de trabajo que estructuren el proceso
        de análisis de licitaciones.
\end{itemize}

Estos objetivos reproducen el alcance registrado del TFG. Su grado de
consecución se evaluará al final de la memoria mediante requisitos, pruebas y
resultados experimentales. En particular, la detección de solvencia y la
comparación entre proveedores no se considerarán validadas únicamente porque
exista una integración técnica con modelos generativos.

\section{Alcance}

El proyecto comprende el desarrollo de una plataforma web orientada a equipos
de consultoras y empresas licitadoras. La solución plantea la gestión de una
cartera de oportunidades mediante pipelines configurables y flujos de trabajo
estructurados. Sobre esta base se integran capacidades de inteligencia
artificial generativa destinadas a apoyar la extracción de información clave,
la evaluación preliminar de solvencia y la generación de documentos de
análisis.

La implementación se organiza como un monorepositorio gestionado con
Turborepo. La interfaz web utiliza Vue 3 y la API se desarrolla con NestJS. El
núcleo funcional comprende la gestión multiempresa de oportunidades,
pipelines y workflows; la configuración de preguntas de control, campos
personalizados y resúmenes; la generación estructurada mediante inteligencia
artificial; y la gestión de documentos PDF utilizados como contexto del
análisis.

El alcance no incluye el descubrimiento ni la importación automática de
licitaciones desde plataformas externas. Tampoco comprende módulos de
organizaciones y contactos, tareas, comentarios, notificaciones, facturación
o exportación de informes a PDF. Estas capacidades se consideran posibles
líneas de trabajo futuro y no se presentan como parte del resultado
implementado.

\section{Estructura de la memoria}

La memoria se organiza en diez capítulos. Tras esta introducción, el
\Cref{chap:planning} presenta la planificación y los costes del proyecto; el
\Cref{chap:requirements} recoge el análisis de requisitos; y el
\Cref{chap:analysis} desarrolla el análisis y diseño del sistema. A
continuación, el \Cref{chap:architecture} describe la arquitectura adoptada y
el \Cref{chap:implementation} detalla la implementación de sus componentes.

El \Cref{chap:testing} expone la estrategia y los resultados de las pruebas y
el \Cref{chap:user-manual} proporciona el manual de usuario de la aplicación.
Finalmente, el \Cref{chap:future-work} analiza las dificultades, limitaciones y
posibles líneas de trabajo futuro, y el \Cref{chap:conclusions} reúne las
conclusiones del proyecto. La declaración sobre el uso de herramientas de
inteligencia artificial generativa como apoyo durante el desarrollo se incluye
en el apéndice~\ref{appendix:ai-assisted-development}.
